% !TeX program = xelatex
\documentclass[12pt,spanish,letterpaper,oneside]{article}
\PassOptionsToPackage{es-noshorthands}{babel}
\def\documenttitle{Plan estratégico integral de participación ciudadana y gobernanza para la cultura de la legalidad en el agro de Nuevo León}
\def\documentsubtitle{Evidencia 3 --- Plan estratégico de solución comunitaria y Estado de derecho (Fase 3)}
\def\documentsubject{Ética y Cultura de la Legalidad --- Ingeniero Agrónomo}
\def\documentauthor{Carlos Mauricio Carvajal Coronado}
\def\coursename{Ética y Cultura de la Legalidad}
\def\coursecode{REF\#261013755115}
\def\universityname{Universidad Autónoma de Nuevo León}
\def\universityfaculty{Facultad de Agronomía}
\def\universitydepartment{Ingeniero Agrónomo}
\def\universitydepartmentimage{departamentos/logos-institucionales.pdf}
\def\universitydepartmentimagecfg{height=2.1cm}
\def\universitylocation{General Escobedo / Marín, Nuevo León, México}
\def\authortable{%
	\begin{Form}
	\begin{tabular}{@{}ll@{}}
		\textbf{Estudiante:} & Carlos Mauricio Carvajal Coronado\\[0.35cm]
		\textbf{Matrícula:} & 1977875\\[0.35cm]
		\textbf{Docente:} & Cátedra de Formación General Universitaria\\[0.35cm]
		\textbf{Unidad de Aprendizaje:} & Ética y Cultura de la Legalidad (Evidencia 3 \textbar\ 5 Pts)\\[0.5cm]
		\multicolumn{2}{l}{Fecha de realización: 14 de octubre de 2026}\\
		\multicolumn{2}{l}{\universitylocation}
	\end{tabular}
	\end{Form}
}

\usepackage{fontspec}
\setmainfont{Arial}
\setsansfont{Arial}
\AddToHook{file/config.tex/after}[aulatex-etica]{%
  \def\documentfontsize{12}
  \def\documentinterline{1}
  \def\documentparindent{0}
  \def\documentparskip{5}
  \def\fontdocument{custom}
  \def\fonttypewriter{custom}
  \def\stylecitereferences{custom}
  \def\portraitstyle{style1}
  \def\hfstyle{style1}
  \def\hfwidthwrap{true}
  \def\hfwidthtitle{0.44}
  \def\hfwidthcourse{0.52}
  \def\predocpageromannumber{false}
  \def\pagemarginleft{2.5}
  \def\pagemarginright{2.5}
  \def\pagemargintop{2.5}
  \def\pagemarginbottom{2.5}
  \def\firstpagemargintop{3.8}
  \def\sectionfontsize{\large}
  \def\sectionspacingtop{8}
  \def\sectionspacingbottom{4}
  \def\ssectionfontsize{\normalsize}
  \def\ssectionspacingtop{6}
  \def\ssectionspacingbottom{3}
  \def\defaultimagefolder{UANL/img/}
  \def\usepdfmetadata{false}
}
% Template:     Informe LaTeX
% Documento:    Núcleo del template
% Versión:      DEV
% Codificación: UTF-8
%
% Autor: Pablo Pizarro R.
%        pablo@ppizarror.com
%
% Manual template: [https://latex.ppizarror.com/informe]
% Licencia MIT:    [https://opensource.org/licenses/MIT]

% -----------------------------------------------------------------------------
% CONFIGURACIONES
% -----------------------------------------------------------------------------
% Definiciones previas
\input{src/defs}

% Ajustes usuario
\input{src/config}

% -----------------------------------------------------------------------------
% IMPORTACIÓN DE LIBRERÍAS
% -----------------------------------------------------------------------------
\input{src/env/imports}

% -----------------------------------------------------------------------------
% IMPORTACIÓN DE FUNCIONES Y ENTORNOS
% -----------------------------------------------------------------------------
\input{src/cmd/core}
\input{src/cmd/math}
\input{src/cmd/equation}
\input{src/cmd/image}
\input{src/cmd/title}
\input{src/cmd/other}
\input{src/cmd/column}
\input{src/env/environments}

% -----------------------------------------------------------------------------
% IMPORTACIÓN DE ESTILOS
% -----------------------------------------------------------------------------
\input{src/style/code}
\input{src/style/other}

% -----------------------------------------------------------------------------
% CONFIGURACIÓN INICIAL DEL DOCUMENTO
% -----------------------------------------------------------------------------
\input{src/cfg/init}

% -----------------------------------------------------------------------------
% PORTADA
% -----------------------------------------------------------------------------
\input{src/page/portrait}
\input{src/page/portrait_config}

% -----------------------------------------------------------------------------
% CONFIGURACIÓN DE PÁGINA Y ENCABEZADOS
% -----------------------------------------------------------------------------
\input{src/cfg/page}

% -----------------------------------------------------------------------------
% TABLA DE CONTENIDOS - ÍNDICE
% -----------------------------------------------------------------------------
\input{src/page/index}

% -----------------------------------------------------------------------------
% CONFIGURACIONES FINALES
% -----------------------------------------------------------------------------
\input{src/cfg/final}

\RemoveFromHook{file/config.tex/after}[aulatex-etica]


\usepackage{tabularx}
\usepackage{booktabs}
\usepackage{amsmath,amssymb}
\usepackage[backend=biber,style=apa,sorting=nyt]{biblatex}
\definecolor{uanlblue}{RGB}{0,55,108}
\definecolor{tableblue}{RGB}{233,240,247}
\def\portraittitlecolor{uanlblue}
\def\sectioncolor{uanlblue}
\def\ssectioncolor{uanlblue}
\def\documenttitlehf{Ética y Legalidad: Evidencia 3 --- Plan Estratégico de Participación Ciudadana}
\hypersetup{colorlinks=true,linkcolor=uanlblue,urlcolor=uanlblue,citecolor=uanlblue,pdfauthor={Carlos Mauricio Carvajal Coronado},pdftitle={Plan estratégico para la cultura de la legalidad en el agro de Nuevo León},pdfsubject={Evidencia 3 - Ética y Legalidad UANL}}
\addbibresource{UANL/ingeniero-agronomo/etica-y-cultura-de-la-legalidad/etica-y-cultura-de-la-legalidad.bib}
\newcolumntype{Y}{>{\raggedright\arraybackslash}X}
\setlength{\emergencystretch}{2.5em}

\begin{document}
\hypersetup{pageanchor=false}
\begingroup
\fontsize{14}{17}\selectfont
\renewcommand{\Huge}{\fontsize{14}{17}\selectfont\bfseries}
\renewcommand{\large}{\fontsize{14}{17}\selectfont}
\templatePortrait
\clearpage
\endgroup
\hypersetup{pageanchor=true}
\templatePagecfg
\singlespacing
\setlength{\tabcolsep}{8pt}
\renewcommand{\arraystretch}{1.18}

\section*{Resumen}
En el presente trabajo se formaliza la Evidencia 3 de la Unidad de Aprendizaje de Ética y Cultura de la Legalidad mediante el diseño de un Plan Estratégico Integral orientado a fortalecer la participación ciudadana, la transparencia y el apego al Estado de derecho en las comunidades rurales del estado de Nuevo León. Con base en el diagnóstico participativo de problemáticas de ilegalidad (acaparamiento ilícito de agua, uso de plaguicidas prohibidos y tala clandestina de mezquitales), el plan estructura cuatro ejes de acción: observatorios agrarios ciudadanos, canal digital seguro de denuncia y rendición de cuentas, certificación ética predial impulsada por la UANL y alianzas interinstitucionales con CONAGUA, PROFEPA y comisarías ejidales. El proyecto integra metas cuantificables, cronograma operativo e indicadores de legalidad verificables.

\vspace{0.4cm}
\noindent\textbf{Palabras clave:} Plan estratégico, participación ciudadana, cultura de la legalidad, gobernanza agraria, observatorios ciudadanos, Estado de derecho.

\section{Introducción y Justificación}
La consolidación de un auténtico Estado constitucional de derecho requiere una ciudadanía activa que vigile, participe y coadyuve en el cumplimiento de las normas que regulan los recursos comunes \parencite{cndh2020}. En el ámbito rural de Nuevo León, la lejanía geográfica de las dependencias fiscalizadoras facilita dinámicas de impunidad que perjudican a pequeños productores y deterioran los ecosistemas \parencite{ley_agraria2022}. Este Plan Estratégico propone una ruta de solución integral fundamentada en la corresponsabilidad cívica y la ética institucional \parencite{cortina2013}.

\section{Ejes Estratégicos del Plan de Intervención}

\subsection{Eje 1: Observatorios Agrarios y Veeduría Comunitaria}
Establecimiento de Comités Ciudadanos de Vigilancia Ambiental y del Agua en los municipios de Marín, General Terán, Cadereyta y Montemorelos. Los comités estarán integrados por ejidatarios, agricultores independientes y estudiantes de la Facultad de Agronomía de la UANL, dotados de capacitación en el marco legal ambiental \parencite{uanl_codigo2021}.

\subsection{Eje 2: Plataforma Tecnológica de Denuncia Protegida y Trazabilidad}
Desarrollo de un canal digital anónimo para georreferenciar tomas clandestinas de agua, descargas industriales no autorizadas en acequias y tiraderos ilegales de envases de agroquímicos. La plataforma generará folios de seguimiento vinculados directamente a la Procuraduría Federal de Protección al Ambiente (PROFEPA) y la Fiscalía General de la República (FGR).

\subsection{Eje 3: Distintivo de Integridad Agraria UANL}
Creación de un programa de acreditación voluntaria para unidades de producción agrícola que certifique el apego estricto a las concesiones de CONAGUA, la adopción del Buen Uso y Manejo de Agroquímicos (BUMA) y el respeto a los derechos laborales de los jornaleros agrícolas conforme a la Constitución \parencite{cpeum2024}.

\section{Matriz Operativa y Cronograma del Plan}

\begin{table}[htbp]
\centering
\small
\caption{Matriz operativa, metas e indicadores del Plan Estratégico de Legalidad.}
\label{tab:plan-legalidad}
\begin{tabularx}{\textwidth}{@{}l Y Y Y@{}}
\toprule
\textbf{Eje de Acción} & \textbf{Acciones Clave} & \textbf{Metas (Año 1--2)} & \textbf{Indicador de Éxito} \\
\midrule
\textbf{Veeduría Ciudadana} & Talleres de derechos agrarios y conformación de 10 comités municipales. & Capacitación de 250 líderes comunitarios y ejidales. & Número de actas de asamblea formalizadas e inscritas ante RAN. \\
\addlinespace
\textbf{Denuncia y Sanción} & Implementación del canal digital seguro y protocolos de verificación. & $100\%$ de reportes canalizados en menos de 72 horas. & Tasa de respuesta ministerial y clausura efectiva de tomas clandestinas. \\
\addlinespace
\textbf{Certificación Ética} & Auditorías agronómicas participativas por brigadas UANL. & Certificación de 40 huertas y predios citrícolas piloto. & Porcentaje de predios con macromedición telemétrica en regla ($>90\%$). \\
\bottomrule
\end{tabularx}
\end{table}

\section{Sustentabilidad y Factibilidad Institucional}
El plan se articula con el Modelo Educativo de la UANL, promoviendo el servicio social comunitario de los estudiantes de Agronomía en vinculación con el Comité Técnico de Aguas Subterráneas (COTAS) y los módulos de los Distritos de Riego 031 y 026. Esta sinergia asegura continuidad operativa, financiamiento mixto y legitimidad social en el territorio.

\section{Conclusiones}
1. La construcción de la legalidad en el campo requiere superar la denuncia pasiva mediante la organización comunitaria estructurada y el empoderamiento legal de los productores agrícolas.
2. La Facultad de Agronomía de la UANL se posiciona como un actor dinamizador de la legalidad y la justicia ambiental en el noreste mexicano, demostrando que la ética aplicada es motor de paz social y prosperidad económica sustentable.

\newpage
\printbibliography[title={Referencias bibliográficas}]
\end{document}

